\section*{अध्याय \ref{ch:b1:counting} --- गणना}

\begin{solution}{exo:b1:counting:1}
स्वतंत्र चरण और गुणन नियम: $26^2 \times 10^3 \times 26^2 = 26^4 \times 1000
= 456\,976\,000$ पट्टिकाएँ। चारों अक्षर जोड़ों में भिन्न होने पर अक्षर वाले
चरण वर्णमाला का $4$-विन्यास बनाते हैं: $26 \times 25 \times 24 \times 23 =
358\,800$ तरीके, अतः $358\,800 \times 1000 = 358\,800\,000$ पट्टिकाएँ।
\end{solution}

\begin{solution}{exo:b1:counting:2}
$\mathrm{ORANGE}$ में $6$ भिन्न अक्षर हैं: $6! = 720$ वर्णविपर्यय।
$\mathrm{BANANA}$ में $6$ अक्षर हैं, जिनमें पुनरावृत्ति है ($3$ a, $2$ n,
$1$ b): प्रत्येक वर्णविपर्यय a के स्थानों ($\binom 63$ चयन), फिर शेष $3$
स्थानों में n के स्थानों ($\binom 32$) से निर्धारित हो जाता है, और b अंतिम
स्थान ले लेता है: $\binom{6}{3}\binom{3}{2} = 20 \times 3 = 60$ वर्णविपर्यय
(समतुल्य रूप से $6!/(3!\,2!\,1!) = 60$)।
\end{solution}

\begin{solution}{exo:b1:counting:3}
कुल: $\binom{12}{4} = 495$। ठीक $2$ स्त्रियाँ: उन्हें चुनिए ($\binom 72 =
21$) और $2$ पुरुष ($\binom 52 = 10$): $210$ समितियाँ। कम से कम एक पुरुष:
``कोई पुरुष नहीं'' का पूरक, $\binom{12}{4} - \binom{7}{4} = 495 - 35 = 460$।
\end{solution}

\begin{solution}{exo:b1:counting:4}
\emph{जन्मदिन:} डिब्बे $12$ माह हैं; $13$ लोगों को $12$ डिब्बों में रखने पर
दो एक ही डिब्बे में आ जाते हैं (\cref{cor:b1:counting:pigeonhole})।

\emph{क्रमागत पूर्णांक:} डिब्बे $n$ युग्म $\{1,2\}, \{3,4\}, \dots, \{2n-1,
2n\}$ हैं, जो $\intint{1}{2n}$ का \omterm{thm:b1:logic:partition}{विभाजन} करते हैं। $n + 1$ पूर्णांक चुनने पर
दो एक ही युग्म में आ जाते हैं, और युग्म के दोनों अवयव क्रमागत हैं।
\end{solution}

\begin{solution}{exo:b1:counting:5}
$1 \leq k \leq n$ के लिए,
\[
k \binom nk = k\,\frac{n!}{k!\,(n-k)!}
= n\,\frac{(n-1)!}{(k-1)!\,(n-k)!} = n \binom{n-1}{k-1}.
\]
$j = k - 1$ से योग करके और पुनःसूचीबद्ध करके:
\[
\sum_{k=0}^{n} k \binom nk = n \sum_{j=0}^{n-1} \binom{n-1}{j}
= n\, 2^{n-1}
\]
\cref{prop:b1:counting:identities} से। (वैकल्पिक: $(1+x)^n = \sum_k \binom
nk x^k$ का अवकलन कीजिए और $x = 1$ रखिए।)
\end{solution}

\begin{solution}{exo:b1:counting:6}
पूर्णतः वर्धमान \omterm{def:b1:logic:map}{प्रतिचित्रण} $f \colon \intint{1}{k} \to \intint{1}{n}$ अपने
प्रतिबिंब से, जो $\intint{1}{n}$ का $k$-उपसमुच्चय है, निर्धारित हो जाता है
(उपसमुच्चय को बढ़ते क्रम में लिखिए); विलोमतः प्रत्येक $k$-उपसमुच्चय ठीक एक
ऐसा \omterm{def:b1:logic:map}{प्रतिचित्रण} देता है। अतः $\binom nk$ पूर्णतः वर्धमान \omterm{def:b1:logic:map}{प्रतिचित्रण} हैं।

यदि $f$ केवल वर्धमान है, तो $g(i) = f(i) + i - 1$ रखिए। तब $g$ पूर्णतः
वर्धमान है (क्रमागत प्रांतिकों के बीच $f$ $\geq 0$ बढ़ता है और $i - 1$ $1$
बढ़ता है) और उसके मान $\intint{1}{n + k - 1}$ में हैं; और $f(i) = g(i) - i +
1$ $\intint{1}{n+k-1}$ में जाने वाले किसी भी पूर्णतः वर्धमान $g$ से $f$ वापस
दे देता है। यह एकैकी आच्छादन है, अतः $\binom{n + k - 1}{k}$ वर्धमान
\omterm{def:b1:logic:map}{प्रतिचित्रण} हैं।
\end{solution}

\begin{solution}{exo:b1:counting:7}
$m + n$ अवयवों वाले \omterm{def:b1:logic:sets}{समुच्चय} $E$ को खंडों $M$ ($m$ अवयव) और $N$ ($n$ अवयव)
में बाँटिए। $E$ का कोई $k$-उपसमुच्चय $M$ के कुछ $j$ अवयव ($0 \leq j \leq k$)
और $N$ के $k - j$ अवयव समेटता है; नियत $j$ के लिए ऐसे $\binom mj
\binom{n}{k-j}$ उपसमुच्चय हैं, और स्थितियाँ $j = 0, \dots, k$
$k$-उपसमुच्चयों का \omterm{thm:b1:logic:partition}{विभाजन} करती हैं। योग नियम वांडरमोंड
सर्वसमिका\index{वांडरमोंड सर्वसमिका} दे देता है।

$m = n = k$ के साथ: $\binom{2n}{n} = \sum_{j=0}^{n} \binom nj \binom{n}{n-j}
= \sum_{j=0}^{n} \binom nj^2$, जहाँ $\binom{n}{n-j} = \binom nj$ का प्रयोग
किया गया।
\end{solution}

\begin{solution}{exo:b1:counting:8}
मान लीजिए $A_d$ $\intint{1}{1000}$ में $d$ के गुणज हैं, अतः $\abs{A_d} =
\lfloor 1000/d \rfloor$। $A_2, A_3, A_5$ के साथ समावेशन--अपवर्जन
(\cref{thm:b1:counting:inclexcl}), यह देखते हुए कि $A_2 \cap A_3 = A_6$
इत्यादि:
\[
500 + 333 + 200 - 166 - 100 - 66 + 33 = 734 .
\]
अतः $734$ पूर्णांक $2$, $3$ या $5$ से विभाज्य हैं।
\end{solution}

\begin{solution}{exo:b1:counting:9}
$2$ अवयवों पर: सभी $2^4 = 16$ \omterm{def:b1:logic:map}{प्रतिचित्रण}, केवल $2$ अचर \omterm{def:b1:logic:map}{प्रतिचित्रण} छोड़कर:
$14$ आच्छादन।

$3$ अवयवों पर: छूट गए मानों पर समावेशन--अपवर्जन से, $4$-समुच्चय से
$3$-समुच्चय में जाने वाले, कम से कम एक मान छोड़ने वाले प्रतिचित्रणों की
संख्या $\binom 31 2^4 - \binom 32 1^4 = 48 - 3 = 45$ है; कुल \omterm{def:b1:logic:map}{प्रतिचित्रण}
$3^4 = 81$; आच्छादन: $81 - 45 = 36$। (जाँच: $4$ से $3$ अवयवों पर कोई आच्छादन
ठीक एक मान को दुहराता है: दुहराया गया मान चुनिए ($3$), उस पर जाने वाला युग्म
($\binom 42 = 6$), और शेष के लिए एक एकैकी आच्छादन ($2$): $3 \times 6 \times
2 = 36$।)
\end{solution}

\begin{solution}{exo:b1:counting:10}
$\N^n$ में $x_1 + \dots + x_n = k$ का हल $k$ तारों और $n - 1$ छड़ों की एक
पंक्ति के रूप में कूटित होता है: $x_1$ तारे लिखिए, एक छड़, $x_2$ तारे, एक
छड़, \dots, और अंत में $x_n$ तारे। यह $k$ तारों और $n - 1$ छड़ों का प्रयोग
करने वाले $k + n - 1$ लंबाई के शब्दों पर एकैकी आच्छादन है, और वे शब्द तारों
के स्थानों से निर्धारित होते हैं: $\binom{n + k - 1}{k}$। पुनरावृत्ति सहित
चयन समीकरण के हलों के अनुरूप हैं ($x_i$ = वस्तु $i$ की प्रतियों की संख्या),
अतः गणना वही है।
\end{solution}

\begin{solution}{exo:b1:counting:11}
$A_i = \{\sigma : \sigma(i) = i\}$ के साथ, $\bigcap_{i \in I} A_i$ का \omterm{def:b1:counting:objects}{क्रमचय}
प्रत्येक $i \in I$ को अचल रखता है और शेष $n - \abs I$ बिंदुओं को स्वतंत्र
रूप से क्रमित करता है: $\abs{\bigcap_{i \in I} A_i} = (n - \abs I)!$।
समावेशन--अपवर्जन:
\[
\Bigl|\bigcup_i A_i\Bigr|
= \sum_{k=1}^{n} (-1)^{k+1} \binom nk (n-k)!
= \sum_{k=1}^{n} (-1)^{k+1} \frac{n!}{k!} ,
\]
क्योंकि $k$ आकार के $I$ उपसमुच्चय $\binom nk$ हैं। अतः
\[
D_n = n! - \Bigl|\bigcup_i A_i\Bigr|
= n!\Bigl(1 - \sum_{k=1}^{n} \frac{(-1)^{k+1}}{k!}\Bigr)
= n! \sum_{k=0}^{n} \frac{(-1)^k}{k!} .
\]
दूसरी सर्वसमिका के लिए: $\intint{1}{n}$ के क्रमचयों $\sigma$ को उनके
अचल-बिंदु \omterm{def:b1:logic:sets}{समुच्चय} $F(\sigma)$ के अनुसार वर्गीकृत कीजिए। नियत $k$-उपसमुच्चय
$F$ के लिए, $F(\sigma) = F$ वाले \omterm{def:b1:counting:objects}{क्रमचय} ठीक पूरक के व्यत्यय हैं: उनकी संख्या
$D_{n-k}$ है। प्रत्येक $k$ के लिए $F$ के $\binom nk$ चयनों पर योग करने से:
$n! = \sum_{k=0}^{n} \binom nk D_{n-k}$।
\end{solution}

\begin{solution}{exo:b1:counting:12}
\emph{पहली सर्वसमिका।} ऐसे युग्म $(A, a)$ गिनिए जिनमें $A \subseteq E$
($\abs E = n$) और $a \in A$। $A$ के आकार के अनुसार: $\sum_k \binom nk k$
युग्म। पहले चिह्नित अवयव चुनने पर: $a$ के लिए $n$ चयन, फिर $A$ पूरा करने के
लिए शेष $n - 1$ अवयवों का कोई भी उपसमुच्चय: $n\,2^{n-1}$ युग्म।

\emph{दूसरी सर्वसमिका।} $a, b \in A$ वाली त्रिकियाँ $(A, a, b)$ गिनिए ($a =
b$ भी संभव है)। आकार के अनुसार: $\sum_k k^2 \binom nk$। सीधे: या तो $a = b$
($n\,2^{n-1}$ त्रिकियाँ, पिछली गणना) या $a \neq b$ ($n(n-1)$ क्रमित चयन, फिर
शेष $n - 2$ अवयवों का कोई भी उपसमुच्चय: $n(n-1)\,2^{n-2}$)। कुल
\[
n\,2^{n-1} + n(n-1)\,2^{n-2} = n\,2^{n-2}\,(2 + n - 1)
= n(n+1)\,2^{n-2} .
\]
\end{solution}

\begin{solution}{pb:b1:counting:1}
\textbf{1.} $D_1 = 0$ (एकमात्र \omterm{def:b1:counting:objects}{क्रमचय} $1$ को अचल रखता है), $D_2 = 1$
(अदला-बदली), $D_3 = 2$ (एक-पंक्ति संकेतन में: $231$ और $312$)। $n = 4$ के
लिए $\sigma(1)$ के अनुसार समूहबद्ध कीजिए: $\sigma(1) = 2$ के साथ व्यत्यय
$2143$, $2341$, $2413$ हैं; $\sigma(1) = 3$ के साथ: $3142$, $3412$, $3421$;
$\sigma(1) = 4$ के साथ: $4123$, $4312$, $4321$। हर समूह में तीन: $D_4 = 9$।

\textbf{2.} ठीक $k$ अचल बिंदुओं वाला \omterm{def:b1:counting:objects}{क्रमचय} अपने अचल-बिंदु \omterm{def:b1:logic:sets}{समुच्चय} $F$ के
चयन ($\binom nk$ तरीके) तथा पूरक पर अपने प्रतिबंधन से निर्धारित होता है, और
वह प्रतिबंधन $n - k$ बिंदुओं का ऐसा \omterm{def:b1:counting:objects}{क्रमचय} होना चाहिए जिसका \emph{कोई} अचल
बिंदु न हो ($D_{n-k}$ तरीके)। दोनों चयन स्वतंत्र हैं और अनुरूपता \omterm{def:b1:logic:inj}{एकैकी
आच्छादक} है: $P_k(n) = \binom nk D_{n-k}$।

\textbf{3.} $P_0(4) = D_4 = 9$; $P_1(4) = \binom41 D_3 = 4 \times 2 = 8$;
$P_2(4) = \binom42 D_2 = 6$; $P_3(4) = \binom43 D_1 = 0$ (तीन अचल बिंदु चौथे
को बाध्य कर देते हैं); $P_4(4) = 1$। योग: $9 + 8 + 6 + 0 + 1 = 24 = 4!$। कोई
मेल नहीं ($9$ स्थितियाँ) ठीक एक मेल ($8$ स्थितियाँ) से --- थोड़े ही अंतर से
--- आगे है।

\textbf{4.} $\sigma(i) = i$ वाले युग्म $(\sigma, i)$ गिनिए। नियत $i$ के लिए
$i$ को अचल रखने वाले \omterm{def:b1:counting:objects}{क्रमचय} शेष $n - 1$ बिंदुओं के \omterm{def:b1:counting:objects}{क्रमचय} हैं: उनकी संख्या
$(n-1)!$ है। अतः युग्मों की संख्या $n \cdot (n-1)! = n!$ है, और यही संख्या
$\sum_\sigma \abs{\mathrm{Fix}(\sigma)}$ भी है। क्रमचयों की संख्या $n!$ से
भाग देने पर: अचल बिंदुओं की औसत संख्या प्रत्येक $n \geq 1$ के लिए ठीक $1$
है।

\textbf{5.} मान लीजिए $\sigma$ $\intint1{n+1}$ का व्यत्यय है और $j =
\sigma(n+1) \in \intint1n$: $n$ संभव मान। \emph{स्थिति $\sigma(j) = n+1$:}
बिंदु $j$ और $n+1$ अदला-बदली कर लेते हैं, और शेष $n - 1$ बिंदुओं पर $\sigma$
का प्रतिबंधन उनका कोई भी व्यत्यय है: $D_{n-1}$ संभावनाएँ। \emph{स्थिति
$\sigma(j) \neq n+1$:} मान लीजिए $i_0 = \sigma^{-1}(n+1)$; यहाँ $i_0 \neq j$
और $i_0 \leq n$। $\intint1n$ पर $\tau$ इस प्रकार परिभाषित कीजिए: $i \neq
i_0$ के लिए $\tau(i) = \sigma(i)$ और $\tau(i_0) = j$। तब $\tau$ $\intint1n$
का \omterm{def:b1:counting:objects}{क्रमचय} है (मान $n+1$ के स्थान पर छूटा हुआ मान $j$ आ गया है), और वह
व्यत्यय है: $\tau(i_0) = j \neq i_0$, तथा अन्यत्र $\tau(i) = \sigma(i) \neq
i$। विलोमतः, $\intint1n$ के व्यत्यय $\tau$ और मान $j$ से $\sigma(n+1) = j$,
$\sigma(\tau^{-1}(j)) = n+1$ रखकर तथा अन्यत्र $\sigma = \tau$ रखकर $\sigma$
वापस मिल जाता है: यह एकैकी आच्छादन है, जो $D_n$ संभावनाएँ देता है। $j$ पर
योग करने से: $D_{n+1} = n(D_n + D_{n-1})$। संख्यात्मक रूप से: $D_5 = 4(9 +
2) = 44$, $D_6 = 5(44 + 9) = 265$।

\textbf{6.} प्रश्न 5 से $D_{n+1} = nD_n + nD_{n-1}$, अतः
\[
u_{n+1} = D_{n+1} - (n+1)D_n = nD_n + nD_{n-1} - (n+1)D_n
= -(D_n - nD_{n-1}) = -u_n .
\]
चूँकि $u_1 = D_1 - 1 \cdot D_0 = -1$, आगमन से $u_n = (-1)^n$ मिलता है,
अर्थात् $n \geq 1$ के लिए $D_n = nD_{n-1} + (-1)^n$।

\textbf{7.} $n$ पर आगमन। आधार: $D_0 = 1 = 0!\,s_0$। पद: $D_{n-1} =
(n-1)!\,s_{n-1}$ मानने पर,
\[
D_n = nD_{n-1} + (-1)^n = n!\,s_{n-1} + (-1)^n
= n!\Bigl(s_{n-1} + \frac{(-1)^n}{n!}\Bigr) = n!\,s_n ,
\]
जो वही सूत्र है। समावेशन--अपवर्जन का प्रयोग कहीं नहीं हुआ: केवल प्रश्न 5 की
संयोजनात्मक पुनरावृत्ति का।

\textbf{8.} त्रिपद पुनर्लेखन, क्रमगुणितों से:
\[
\binom nk \binom kj
= \frac{n!}{k!\,(n-k)!} \cdot \frac{k!}{j!\,(k-j)!}
= \frac{n!}{j!\,(n-j)!} \cdot \frac{(n-j)!}{(k-j)!\,(n-k)!}
= \binom nj \binom{n-j}{k-j} .
\]
अब $a_k = \sum_j \binom kj b_j$ प्रतिस्थापित कीजिए और दोनों \omterm{def:b1:counting:card}{परिमित} योगों की
अदला-बदली कीजिए:
\[
\sum_{k=0}^{n} (-1)^{n-k} \binom nk a_k
= \sum_{j=0}^{n} b_j \binom nj
\sum_{k=j}^{n} (-1)^{n-k} \binom{n-j}{k-j}
= \sum_{j=0}^{n} b_j \binom nj
\sum_{i=0}^{n-j} (-1)^{(n-j)-i} \binom{n-j}{i} .
\]
भीतरी योग $(1 + (-1))^{n-j} = 0^{n-j}$ का प्रसार है (द्विपद प्रमेय,
\cref{thm:b1:counting:binomial}): वह $j < n$ के लिए लुप्त हो जाता है और $j =
n$ के लिए $1$ के बराबर है। केवल $j = n$ बचता है, और दायाँ पक्ष $b_n$ है,
जैसा दावा किया गया था।

\textbf{9.} सममिति $\binom nk = \binom n{n-k}$ से \cref{exo:b1:counting:11}
की सर्वसमिका $n! = \sum_{k=0}^n \binom nk D_k$ के रूप में फिर से लिखी जाती
है। $a_n = n!$ और $b_k = D_k$ के साथ प्रश्न 8 लगाइए:
\[
D_n = \sum_{k=0}^{n} (-1)^{n-k} \binom nk k!
= \sum_{k=0}^{n} (-1)^{n-k} \frac{n!}{(n-k)!}
= n! \sum_{j=0}^{n} \frac{(-1)^j}{j!} ,
\]
$j = n - k$ से पुनःसूचीबद्ध करने पर: वही सूत्र तीसरी बार।

\textbf{10.} $D_n = n!\,s_n$ (प्रश्न 7), अतः
\[
\Bigl| D_n - \frac{n!}{\eu} \Bigr|
= n!\,\abs{s_n - \eu^{-1}} < \frac{n!}{(n+1)!} = \frac1{n+1} .
\]

\textbf{11.} $n \geq 1$ के लिए $\frac1{n+1} \leq \frac12$, और प्रश्न 10 की
असमिका कठोर है: $D_n$ $n!/\eu$ से $< \frac12$ दूरी पर है, अतः वही अद्वितीय
निकटतम पूर्णांक है। $n = 0$ के लिए परिबंध केवल $< 1$ दूरी देता है, और सचमुच
दावा वहाँ विफल हो जाता है: $0!/\eu \approx 0.368$ का निकटतम पूर्णांक $0$ है,
जबकि $D_0 = 1$।

\textbf{12.} $\eu^{-1} - s_n = \sum_{k \geq n+1} (-1)^k/k!$ कठोरतः घटते पदों
वाली एकांतरित श्रेणी है, अतः उसका चिह्न उसके पहले पद $(-1)^{n+1}/(n+1)!$ का
चिह्न है। अतः $s_n - \eu^{-1}$ का चिह्न $(-1)^n$ का चिह्न है: $n$ सम होने पर
$s_n > \eu^{-1}$ और $D_n = n!\,s_n > n!/\eu$; $n$ विषम होने पर $D_n <
n!/\eu$।

\textbf{13.} $D_7 = 6(265 + 44) = 6 \times 309 = 1854$; $D_8 = 7(1854 + 265)
= 7 \times 2119 = 14\,833$; $D_9 = 8(14\,833 + 1854) = 8 \times 16\,687 =
133\,496$; $D_{10} = 9(133\,496 + 14\,833) = 9 \times 148\,329 =
1\,334\,961$। जाँच: $10!/\eu = 3\,628\,800 / 2.718281828 \approx
1\,334\,960.92$, जिसका निकटतम पूर्णांक $1\,334\,961$ है --- और $D_{10} >
10!/\eu$, जैसा प्रश्न 12 सम $n$ के लिए बताता है।

\textbf{14.} $\abs{p_n - \eu^{-1}} = \abs{s_n - \eu^{-1}} < \frac1{(n+1)!}$।
$n = 6$ के लिए: $p_6 = 265/720 = 0.36806$ (पाँच दशमलव), जबकि $\eu^{-1} =
0.36788$; अंतर $1/7! = 1/5040 < 2 \times 10^{-4}$ से कम है। परिबंध
$1/(n+1)!$ इतनी तेज़ी से ढहता है कि दर्जन भर पत्रों पर ही प्रायिकता अनेक
दशमलव स्थानों तक स्थिर हो जाती है: ``लगभग $36.8\%$'' वाला उत्तर हर
व्यावहारिक प्रयोजन के लिए $n$ से स्वतंत्र है --- यही इस समस्या का प्रसिद्ध
विस्मय है।

\textbf{15.} प्रश्न 2 और $D_m = m!\,s_m$ से:
\[
\frac{P_k(n)}{n!} = \frac{\binom nk D_{n-k}}{n!}
= \frac{D_{n-k}}{k!\,(n-k)!} = \frac{s_{n-k}}{k!}
\;\longrightarrow\; \frac{\eu^{-1}}{k!}
\]
जब $k$ नियत रखते हुए $n \to \infty$, क्योंकि $s_{n-k} \to \eu^{-1}$। सीमांत
मान $\eu^{-1}/k!$ ($k \in \N$) प्राचल $1$ के प्वासों बंटन के भार हैं।

\textbf{16.} $i \neq j$, $\sigma(i) = i$, $\sigma(j) = j$ वाली त्रिकियाँ
$(\sigma, i, j)$ गिनिए। पहले क्रमित युग्म चुनने पर: $n(n-1)$ तरीके; $i$ और
$j$ दोनों को अचल रखने वाले \omterm{def:b1:counting:objects}{क्रमचय} शेष $n - 2$ बिंदुओं के \omterm{def:b1:counting:objects}{क्रमचय} हैं: उनकी
संख्या $(n-2)!$ है। कुल: $n(n-1)(n-2)! = n!$। इसके बदले पहले $\sigma$ पर योग
करने से प्रत्येक $\sigma$ के लिए भिन्न अचल बिंदुओं के क्रमित युग्म गिने जाते
हैं: $\abs{\mathrm{Fix}(\sigma)}(\abs{\mathrm{Fix}(\sigma)}-1)$। अतः वही
सर्वसमिका मिलती है; $n!$ से भाग देने पर
$\abs{\mathrm{Fix}}(\abs{\mathrm{Fix}} - 1)$ का औसत $1$ है, अतः
$\abs{\mathrm{Fix}}^2$ का औसत $1 + 1 = 2$ है और प्रसरण $2 - 1^2 = 1$ है।

\textbf{17.} अनुपात $1 - p_n$: $n = 4$ के लिए $1 - \frac 9{24} =
\frac{15}{24} = 0.6250$; $n = 5$ के लिए $1 - \frac{44}{120} = \frac{76}{120}
= 0.6333$; $n = 6$ के लिए $1 - \frac{265}{720} = \frac{455}{720} = 0.6319$।
सभी $1 - \eu^{-1} \approx 0.6321$ के एक प्रतिशत के भीतर, उसके चारों ओर दोलन
करते हुए।

\textbf{18.} सीधे:
\[
s_{n+2} - s_n = \frac{(-1)^{n+1}}{(n+1)!} +
\frac{(-1)^{n+2}}{(n+2)!}
= (-1)^{n+1}\Bigl(\frac1{(n+1)!} - \frac1{(n+2)!}\Bigr),
\]
और कोष्ठक $> 0$ है। $n$ सम होने पर अंतर ऋणात्मक है: $s_{n+2} < s_n$, अतः
$p_0 > p_2 > p_4 > \dots$; $n$ विषम होने पर वह धनात्मक है: $p_1 < p_3 < p_5
< \dots$। प्रश्न 12 (सम मान $\eu^{-1}$ से ऊपर, विषम नीचे) और प्रश्न 14
($\eu^{-1}$ से दूरी $0$ की ओर जाती है) के साथ मिलाकर: दोनों सीढ़ियाँ
$\eu^{-1}$ को अपने बीच दबा लेती हैं।

\textbf{19.} एक पूर्ण निष्कर्षण समरूप यादृच्छिक \omterm{def:b1:counting:objects}{क्रमचय} है, जो व्यत्यय होने
पर वैध है: प्रायिकता $p_n \approx \eu^{-1}$। उद्धृत तथ्य से सफलता तक
निष्कर्षणों की औसत संख्या $1/p_n$ है, और प्रश्न 14 से $1/p_n \approx \eu$
मिलता है, जिसकी त्रुटि छोटे $n$ पर भी नगण्य है। अतः पुनरारंभ वाला गुप्त
उपहार औसतन लगभग $\eu \approx 2.72$ पूर्ण निष्कर्षणों की पड़ती है --- चाहे
कार्यालय में $6$ लोग हों या $600$।

\textbf{20.} $j \geq 2$ नियत कीजिए और मान $\sigma(1) = j$ पर प्रश्न 5 का
वर्गीकरण चलाइए। यदि $\sigma(j) = 1$: शेष $n - 2$ बिंदु कोई भी व्यत्यय धारण
करते हैं, $D_{n-2}$ तरीके। यदि $\sigma(j) \neq 1$: प्रश्न 5 की भाँति ही
\omterm{def:b1:logic:map}{पूर्वप्रतिबिंब} $i_0 = \sigma^{-1}(1)$ की दिशा $j$ की ओर मोड़ दीजिए; यह $n -
1$ बिंदुओं $\{2, \dots, n\}$ के व्यत्ययों के साथ एकैकी आच्छादन है: $D_{n-1}$
तरीके। कुल $D_{n-1} + D_{n-2}$, जो प्रत्येक $j$ के लिए वही है। $j$ के $n -
1$ मानों पर योग करने से: $D_n = (n-1)(D_{n-1} + D_{n-2})$, जो गुणनखंड $n -
1$ प्रकट कर देता है: $(n-1) \mid D_n$।

\textbf{21.} दावा: $D_n$ विषम है यदि और केवल यदि $n$ सम हो। $D_n = nD_{n-1}
+ (-1)^n$, अर्थात् $D_n \equiv nD_{n-1} + 1 \pmod 2$, का प्रयोग करते हुए
आगमन। आधार: $D_1 = 0$ सम है, $n = 1$ विषम: दावा सत्य है। यदि $n$ सम है, तो
$nD_{n-1}$ सम है और $D_n \equiv 1$: विषम, जैसा दावा किया गया। यदि $n$ विषम
है, तो $n - 1$ सम है, अतः परिकल्पना से $D_{n-1}$ विषम है, और $D_n \equiv
D_{n-1} + 1 \equiv 0$: सम। आगमन बंद हो जाता है।

\textbf{22.} $D_n = nD_{n-1} + (-1)^n$ को $n$ के सापेक्ष घटाने पर पहला पद
समाप्त हो जाता है: $D_n \equiv (-1)^n \pmod n$। $n = 10$ के लिए: $(-1)^{10}
= 1$, और सचमुच $D_{10} = 1\,334\,961$ का अंतिम अंक $1$ है।

\textbf{23.} $n \geq 3$ के लिए $D_{n-1} \geq 1$, और प्रश्न 6 की पुनरावृत्ति
को $D_{n-1}$ से भाग देने पर $D_n/D_{n-1} = n + (-1)^n/D_{n-1}$ मिलता है,
जहाँ $\abs{(-1)^n/D_{n-1}} \leq 1$ है और वह तेज़ी से $0$ की ओर जाता है।
संगति: यदि $D_n \approx n!/\eu$, तो $D_n/D_{n-1} \approx n!/(n-1)! = n$ ---
अनुपात में गुणनखंड $\eu$ कट जाता है, और पुनरावृत्ति उसकी $1/D_{n-1}$
यथार्थता तक पुष्टि कर देती है।

\textbf{24.} (क) गुणन और योग नियम हर गणना के मूल में हैं: प्रश्न 2 और 5
क्रमचयों के समुच्चयों को स्वतंत्र चरणों में विभाजित करते हैं। (ख) दुहरी गणना
ने $D_n$ के लिए कोई सूत्र लिए बिना ही अचल बिंदुओं की संख्या का माध्य (प्रश्न
4) और प्रसरण (प्रश्न 16) दे दिया। (ग) द्विपद प्रमेय ने वह एकांतरित भीतरी योग
$(1-1)^{n-j}$ आँका जिससे द्विपद प्रतिलोमन काम करता है (प्रश्न 8)। (घ)
एकांतरित-श्रेणी परिबंध ने यथार्थ किंतु अपारदर्शी योग $n!\,s_n$ को ``$n!/\eu$
का निकटतम पूर्णांक'' जैसे पारदर्शी \omterm{def:b1:logic:statement}{कथन} में बदल दिया (प्रश्न 10--14)।

\textbf{25.} समावेशन--अपवर्जन (\cref{ex:b1:counting:derangement} और
\cref{exo:b1:counting:11}) वैचारिक उपपत्ति है: वह एकांतरित योग को अधिगणना के
सुधारों के रूप में समझाती है, और ``बुरे'' समुच्चयों के किसी भी परिवार से
बचने वाले अवयवों की गणना तक अक्षरशः सामान्यीकृत हो जाती है। पुनरावृत्ति वाला
रास्ता (प्रश्न 5--7) सबसे तेज़ संगणना करता है --- रैखिक समय, यथार्थ पूर्णांक
अंकगणित, कोई क्रमगुणित नहीं --- और भाग V के अंकगणितीय तथ्यों का स्रोत है।
द्विपद प्रतिलोमन (प्रश्न 8--9) सूत्र को एक व्यापक रूपांतर के भीतर रख देता
है, जो वहाँ-वहाँ फिर प्रकट होगा जहाँ सर्वसमिकाओं के दो त्रिभुजीय निकाय
आमने-सामने आते हैं। और $\eu$ का प्रकट होना स्वयं सूत्र से सबसे अच्छा समझ में
आता है: व्यत्ययों का अनुपात $\eu^{-1}$ की श्रेणी का आंशिक योग $s_n$ है, अतः
मोंमोर के लिफ़ाफ़े, ऑयलर के संकेतन से तीन दशक पहले ही, संख्या $\eu$ की
संगणना कर रहे थे।
\end{solution}
